\uuid{J2WS}
\titre{Logique}
\auteur{bodin}
\organisation{exo7}
\licence{CC BY-NC-SA 4.0}
\datecreate{2020-01-21}
\niveau{L1}
\module{Algèbre}
\chapitre{Logique, ensemble, raisonnement}
\sousChapitre{Autre}
\difficulte{}
\theme{Logique}
\exo7qcmid{arnaud.arnaud.009}

\contenu{
\begin{question}
\qclassification{100.99}
On lance de façon aléatoire deux dés identiques à $6$
faces (numérotées de $1$ à $6$). On ne tient pas compte de
l'ordre, par exemple le tirage $1$ puis $5$ est le même que $5$
puis $1$, mais les tirages $3$ puis $3$, et $3$ puis $4$ sont
distincts.
\begin{answers}
    \bad{Il y a $36$ tirages distincts possibles.}

    \good{Il y a $30$ tirages distincts possibles.}

    \bad{Il y a $21$ tirages distincts possibles.}

    \good{La somme des deux chiffres a strictement plus de chances d'être $7$ que $2$.}

    \bad{La somme des deux chiffres a strictement plus de chances d'être $\ge 11$ que $\le 3$.}
\end{answers}
\end{question}
}
